
\subsubsection{7.1 Summary}

Experiments on Waterbirds with ResNet-50 yielded a null result---spurious and core alignments were both approximately 0.05, statistically indistinguishable. Investigation traced this to measurement apparatus limitation: accumulating one gradient per epoch produced a 10-dimensional subspace in a 25.6-million-parameter space.

Contributions are methodological:

\begin{enumerate}
\item Minimum requirements for gradient subspace analysis are established: multi-batch accumulation within early epochs.
\item A reproducible failure case is documented that future gradient-based methods should avoid.
\item Guidance is provided for valid testing of gradient orthogonalization hypotheses.
\end{enumerate}

\subsubsection{7.2 Future Directions}

\textbf{Immediate:} Re-implement gradient accumulation using streaming SVD across all batches during epochs 1-10.

\textbf{Medium-term:} If valid measurement confirms spurious-core separation, implement full PGO mechanism and evaluate against baselines.

\textbf{Long-term:} Extend to other architectures and domains where gradient subspace properties may differ.

\subsubsection{7.3 Closing}

The question of whether early gradients capture spurious directions remains open. What has been established is what valid measurement requires. Future work should validate accumulation density before claiming directional findings.

